\section{Prospective confirmations beyond order 36}
\label{sec:prospective}

The double-ladder formulas were tested prospectively before their general
proofs were completed.  A timestamped, read-only prediction record froze the
following values before any of the three target graphs was constructed or
analyzed:
\[
\begin{array}{c|rr}
G & |\Hcal(G)| & \hsep(G)\\ \hline
D(9,9) & 86 & 60\\
D(11,9) & 104 & 71\\
D(11,11) & 126 & 84.
\end{array}
\]

\begin{computationalresult}[Prediction-locked double-ladder test]
\label{res:prospective-ladders}
All six frozen values were subsequently confirmed exactly.  Each graph was
obtained from the certified \(D(9,7)\) by the specified facial square
insertions.  For each target, a path-based enumerator and an independent
perfect-matching/complement enumerator produced identical complete normalized
Hamiltonian-cycle universes.  Explicit primal and matching packing
certificates established the three displayed \(\hsep\) values, and the
standard-library verifier passed every graph, cycle, ordered edge-pair, lower
certificate, and manifest check.
\end{computationalresult}

Historically, these were genuinely prospective checks of both formulas.
Mathematically, the Hamiltonian-cycle values now also follow from
\cref{thm:cycle-formula}, and the \(\hsep\) values from
\cref{thm:exact-double-ladder}.  The computations remain useful independent
finite checks of the symbolic descriptions and of three square-insertion
lifts.

No complete Barnette census was run at order 40, 44, or 48 for this test.
Accordingly, Computational Result~\ref{res:prospective-ladders} proves neither
\(\MB(40)=60\) nor any value of \(\MB(44)\) or \(\MB(48)\); it makes no
order-level extremal claim.
